\chapter{กรณีศึกษาภาคสนาม การประยุกต์ใช้เทคโนโลยี AIoT ในสวนทุเรียนจังหวัดจันทบุรี}
\label{chap:ch08}

\section*{แผนบริหารการสอนประจำบทที่ 8}
\addcontentsline{toc}{section}{แผนบริหารการสอนประจำบทที่ 8}

\noindent\textbf{หัวข้อเนื้อหาประจำบท}
\begin{enumerate}
    \item บริบทแปลงทดสอบจริงและผลลัพธ์จากโครงการสวนพ่อรวย จันทบุรี
    \item โครงการทุเรียนประณีตวิทยาคมและการถ่ายทอดสู่เยาวชนและชุมชน
    \item การวิเคราะห์ความคุ้มค่าทางการเงินและผลตอบแทนการลงทุน
\end{enumerate}

\noindent\textbf{วัตถุประสงค์เชิงพฤติกรรม}
\begin{enumerate}
    \item อธิบายหลักการ ทฤษฎี และสถาปัตยกรรมเทคโนโลยีประจำบทที่ 8 ได้อย่างถูกต้อง
    \item ประยุกต์ใช้เครื่องมือ อุปกรณ์เซนเซอร์ และระบบปัญญาประดิษฐ์เพื่อการจัดการสวนทุเรียนได้อย่างมีประสิทธิภาพ
    \item วิเคราะห์ สรุปผลข้อมูล และแก้ไขปัญหาเชิงฟิสิกส์เกษตรและคุณภาพดินได้อย่างเป็นระบบ
\end{enumerate}

\noindent\textbf{การวัดและประเมินผล}
\begin{table}[H]
\centering\small
\begin{tabularx}{\textwidth}{p{0.55\textwidth} p{0.25\textwidth} c}
\toprule
\textbf{เกณฑ์การประเมินผล} & \textbf{เครื่องมือวัดผล} & \textbf{สัดส่วนคะแนน} \\
\midrule
1. ทักษะการปฏิบัติงานและการใช้อุปกรณ์ & แบบประเมินทักษะภาคสนาม & 40\% \\
2. รายงานข้อมูลและการวิเคราะห์ผล & การตรวจสมุดบันทึกและดาต้าชีท & 30\% \\
3. แบบฝึกหัดและคำถามท้ายบท & แบบทดสอบท้ายบทเรียน & 20\% \\
4. การมีส่วนร่วมและความรับผิดชอบ & การเข้าเรียนและการตรงต่อเวลา & 10\% \\
\midrule
\textbf{รวมทั้งสิ้น} & & \textbf{100\%} \\
\bottomrule
\end{tabularx}
\end{table}
\newpage

% ==========================================================
% เนื้อหาบทเรียนประจำบทที่ 8
% ==========================================================

\section{บริบทแปลงทดสอบจริงและผลลัพธ์จากโครงการสวนพ่อรวย จันทบุรี}
\index{บริบทแปลงทดสอบจริงและผลลัพธ์จากโครงการสวนพ่อรวย จันทบุรี}

การศึกษาและทำความเข้าใจเกี่ยวกับบริบทแปลงทดสอบจริงและผลลัพธ์จากโครงการสวนพ่อรวย จันทบุรี (Suan Por Ruay Commercial Orchard Case Study) มีความสำคัญอย่างยิ่งในระบบเกษตรอัจฉริยะ โดยมีรายละเอียดสาระสำคัญดังนี้

สวนพ่อรวย ตั้งอยู่ในพื้นที่ตำบลชากไทย อำเภอเขาคิชฌกูฏ จังหวัดจันทบุรี เป็นแปลงปลูกทุเรียนพันธุ์หมอนทองเชิงพาณิชย์ขนาด 40 ไร่ ซึ่งได้รับเลือกเป็นแปลงทดสอบระบบ AIoT และการให้น้ำแม่นยำสูง

สภาวะก่อนการปรับปรุงเทคโนโลยี:
- สวนใช้วิธีการให้น้ำโดยกะเวลาตามความรู้สึกของแรงงาน ทำให้บางบริเวณได้รับน้ำมากเกินไปจนดินแฉะและเกิดโรครากเน่า ในขณะที่บริเวณเนินสูงขาดแคลนน้ำ
- ต้นทุนค่าน้ำมันและไฟฟ้าสำหรับสูบน้ำเฉลี่ยสูงถึง 18,000 บาทต่อเดือน
- มีปัญหาผลแตกและอาการเต่าเผาในฤดูร้อนเนื่องจากความชื้นผันผวน

ผลลัพธ์หลังการติดตั้งระบบ AIoT และเซนเซอร์ควบคุมอัตโนมัติ:
1. การใช้น้ำลดลง 31.4\% จากปริมาณการใช้น้ำเฉลี่ย 12,000 ลบ.ม./เดือน เหลือ 8,230 ลบ.ม./เดือน
2. ค่าใช้จ่ายด้านพลังงานลดลง 28.5\%
3. สัดส่วนผลผลิตทุเรียนเกรดส่งออก (AB Grade) เพิ่มขึ้นจาก 68\% เป็น 84\%
4. ลดการสูญเสียผลร่วงหล่นจากพายุฤดูร้อนด้วยการแจ้งเตือนสภาพอากาศล่วงหน้า

\begin{casebox}[บทเรียนความสำเร็จจากการลงพื้นที่จริง]
ความสำเร็จของการทำเกษตรอัจฉริยะมิได้ขึ้นอยู่กับความซับซ้อนของเทคโนโลยีเพียงอย่างเดียว หากแต่อยู่ที่ความเข้าใจบริบทของเกษตรกร ความง่ายในการใช้งาน และความทนทานต่อสภาพแวดล้อมเขตร้อนชื้น
\end{casebox}

\section{โครงการทุเรียนประณีตวิทยาคมและการถ่ายทอดสู่เยาวชนและชุมชน}
\index{โครงการทุเรียนประณีตวิทยาคมและการถ่ายทอดสู่เยาวชนและชุมชน}

การศึกษาและทำความเข้าใจเกี่ยวกับโครงการทุเรียนประณีตวิทยาคมและการถ่ายทอดสู่เยาวชนและชุมชน (Praneet Witthayakhom Model School Project) มีความสำคัญอย่างยิ่งในระบบเกษตรอัจฉริยะ โดยมีรายละเอียดสาระสำคัญดังนี้

โครงการขยายผลเทคโนโลยี AIoT สู่สถานศึกษา ณ โรงเรียนประณีตวิทยาคม อำเภอเขาสมิง จังหวัดตราด ภายใต้การสนับสนุนงบประมาณวิจัยและบริการวิชาการของมหาวิทยาลัยราชภัฏรำไพพรรณี

วัตถุประสงค์และกิจกรรมหลัก:
- จัดสร้างแปลงสาธิตทุเรียนอัจฉริยะ (Smart Durian Learning Lab) ขนาด 5 ไร่ในโรงเรียน
- พัฒนาหลักสูตรท้องถิ่น "นวัตกรน้อยเกษตรอัจฉริยะ" สำหรับนักเรียนระดับมัธยมศึกษา
- ฝึกอบรมนักเรียนและชุมชนให้สามารถประกอบ ติดตั้ง และซ่อมบำรุงชุดเซนเซอร์ LoRaWAN และระบบน้ำอัตโนมัติได้ด้วยตนเอง
- เกิดการถ่ายทอดองค์ความรู้จากนักเรียนสู่แปลงสวนของครอบครัวในชุมชน เกิดความยั่งยืนทางเศรษฐกิจและสังคม

\begin{casebox}[บทเรียนความสำเร็จจากการลงพื้นที่จริง]
ความสำเร็จของการทำเกษตรอัจฉริยะมิได้ขึ้นอยู่กับความซับซ้อนของเทคโนโลยีเพียงอย่างเดียว หากแต่อยู่ที่ความเข้าใจบริบทของเกษตรกร ความง่ายในการใช้งาน และความทนทานต่อสภาพแวดล้อมเขตร้อนชื้น
\end{casebox}

\section{การวิเคราะห์ความคุ้มค่าทางการเงินและผลตอบแทนการลงทุน}
\index{การวิเคราะห์ความคุ้มค่าทางการเงินและผลตอบแทนการลงทุน}

การศึกษาและทำความเข้าใจเกี่ยวกับการวิเคราะห์ความคุ้มค่าทางการเงินและผลตอบแทนการลงทุน (Economic Feasibility and Return on Investment) มีความสำคัญอย่างยิ่งในระบบเกษตรอัจฉริยะ โดยมีรายละเอียดสาระสำคัญดังนี้

การวิเคราะห์ผลตอบแทนการลงทุน (Return on Investment; ROI) ของการติดตั้งระบบ AIoT ในสวนทุเรียนขนาด 20 ไร่ (ประมาณ 400 ต้น):

1. ต้นทุนการลงทุนเริ่มต้น (Capital Expenditure; CapEx):
   - สถานีตรวจอากาศ AIoT และเกตเวย์ LoRaWAN: 25,000 บาท
   - โพรบวัดดินและโหนดเซนเซอร์ 4 จุด: 32,000 บาท
   - ชุดควบคุมวาล์วไฟฟ้าและท่อประปาแยกโซน: 38,000 บาท
   - รวมเงินลงทุนเริ่มต้น: 95,000 บาท
2. ผลประโยชน์ที่ได้รับต่อปี (Operational Benefits):
   - ประหยัดค่าน้ำมันและค่าไฟฟ้าสูบน้ำ: 24,000 บาท/ปี
   - ลดความเสียหายของผลผลิตทุเรียนจากโรครากเน่าและผลแตก: 80,000 บาท/ปี
   - รายได้เพิ่มขึ้นจากการได้ทุเรียนเกรดพรีเมียม: 120,000 บาท/ปี
   - รวมผลประโยชน์สุทธิต่อปี: 224,000 บาท/ปี
3. ระยะเวลาคืนทุน (Payback Period): ต่ำกว่า 6 เดือน ซึ่งยืนยันความคุ้มค่าเชิงพาณิชย์อย่างชัดเจน

\section{ขั้นตอนและวิธีปฏิบัติการทดลอง}
\index{ขั้นตอนการทดลองประจำบทที่ 8}

\subsection{ขั้นตอนการออกแบบผังการวางระบบ AIoT สำหรับสวนทุเรียนขนาด 20 ไร่}
\begin{sensorbox}[ขั้นตอนการออกแบบผังการวางระบบ AIoT สำหรับสวนทุเรียนขนาด 20 ไร่]
1. วัดพิกัดแปลงด้วยเครื่องรับสัญญาณดาวเทียม GNSS นำเข้าโปรแกรม QGIS
2. แบ่งพื้นที่ออกเป็น 4 โซนย่อยตามลักษณะความลาดชันและชนิดดิน
3. กำหนดตำแหน่งเกตเวย์กลางและเสาส่งสัญญาณบนหลังคาบ้านพักเจ้าของสวน
4. วางแผนแนวท่อเมนและตำแหน่งติดตั้งกล่องวาล์วไฟฟ้าประจำแต่ละโซน
5. คำนวณงบประมาณรายการวัสดุและอุปกรณ์ (Bill of Quantities; BOQ)
\end{sensorbox}

\subsection{การฝึกอบรมเกษตรกรและนักเรียนในการอ่านค่าแดชบอร์ดและการบำรุงรักษา}
\begin{sensorbox}[การฝึกอบรมเกษตรกรและนักเรียนในการอ่านค่าแดชบอร์ดและการบำรุงรักษา]
1. จัดการฝึกอบรมเชิงปฏิบัติการติดตั้งแอปพลิเคชันบนสมาร์ทโฟนของเกษตรกร
2. สอนการแปลความหมายของกราฟความชื้นดิน: จุดเตือนให้น้ำ (Wilting point) และจุดอิ่มตัว (Saturation)
3. สาธิตการทำความสะอาดแผงโซลาร์เซลล์และการตรวจเช็กขั้วต่อแบตเตอรี่ทุก 3 เดือน
4. ฝึกหัดการสั่งการเปิด-ปิดน้ำด้วยตนเองและการแก้ปัญหาเบื้องต้นเมื่อเซนเซอร์หยุดส่งข้อมูล
\end{sensorbox}

\subsection{การเก็บข้อมูลและประเมินผลสัมฤทธิ์ก่อนและหลังการใช้เทคโนโลยี}
\begin{sensorbox}[การเก็บข้อมูลและประเมินผลสัมฤทธิ์ก่อนและหลังการใช้เทคโนโลยี]
1. บันทึกปริมาณการใช้น้ำ (ลูกบาศก์เมตร) จากมาตรวัดน้ำทุกสัปดาห์
2. บันทึกค่าใช้จ่ายพลังงานและจำนวนชั่วโมงการทำงานของปั๊มน้ำ
3. ตรวจนับจำนวนผลทุเรียนต่อต้น และสัดส่วนเกรดผลผลิต ณ จุดรับซื้อ
4. สรุปรายงานเปรียบเทียบเชิงเศรษฐศาสตร์และนำเสนอผลต่อกลุ่มวิสาหกิจชุมชน
\end{sensorbox}

\section{ตารางบันทึกผลการทดลองและการอภิปรายผล}
\begin{table}[H]
\centering\small
\caption{ตารางบันทึกผลการตรวจวัดและข้อมูลพารามิเตอร์ประจำบทที่ 8}
\label{tab:ch08_datasheet}
\begin{tabularx}{\textwidth}{p{0.3\textwidth} p{0.3\textwidth} X}
\toprule
\textbf{พารามิเตอร์ / การทดสอบ} & \textbf{ค่าที่ตรวจวัดได้} & \textbf{การแปลผลและการวิเคราะห์} \\
\midrule
จุดตรวจวัดที่ 1 (บริเวณดินดอน) & ค่าความชื้นอยู่ในเกณฑ์ปกติ & ปริมาณน้ำเพียงพอต่อการเจริญ \\
จุดตรวจวัดที่ 2 (บริเวณที่ลุ่ม) & มีการสะสมความชื้นสูง & ควรชะลอการให้น้ำเพื่อป้องกันรากเน่า \\
สถานะสัญญาณโครงข่าย & RSSI: -85 dBm, SNR: +8 dB & คุณภาพสัญญาณ LoRaWAN ยอดเยี่ยม \\
\bottomrule
\end{tabularx}
\end{table}

\section{คำถามทบทวนและแบบฝึกหัดท้ายบท}
\begin{enumerate}
    \item จากกรณีศึกษาสวนพ่อรวย ปัจจัยใดส่งผลให้สัดส่วนทุเรียนเกรดส่งออกเพิ่มขึ้นจาก 68\% เป็น 84\%
    \item จงคำนวณระยะเวลาคืนทุน (Payback Period) ของโครงการ AIoT หากเงินลงทุนเริ่มต้นเท่ากับ 120,000 บาท และสร้างผลประโยชน์สุทธิได้ 180,000 บาทต่อปี
    \item การขยายผลเทคโนโลยี AIoT สู่เยาวชนในโรงเรียนประณีตวิทยาคมมีคุณค่าเชิงยุทธศาสตร์ต่อการพัฒนาชุมชนเกษตรกรรมอย่างไร
\end{enumerate}
